\subsection{Dynamic Typing in Assignment}

Python is dynamically typed in the sense that names do not have fixed declared types. The object has a type through \texttt{ob\_type}; the name is only a binding to an object.

\subsubsection{Names Do Not Have Fixed Declared Types}

In C, \texttt{int score;} fixes the storage type for the lifetime of that declaration. In Python, \texttt{score = 42} binds \texttt{score} to a \texttt{PyLongObject}. Rebinding \texttt{score = "text"} replaces the binding with a string object. The name carries no type annotation unless the programmer adds one (and even annotations do not enforce runtime behavior by default).

The practical rule:

\begin{quote}
Do not ask what type a Python name has permanently. Ask what object the name currently references.
\end{quote}

Available methods come from the current object's type, not from the spelling of the name. After \texttt{score = 42}, \texttt{bit\_length} exists; after \texttt{score = "text"}, \texttt{upper} exists instead.

\subsubsection{Rebinding the Same Name to Objects of Different Types}

The same name can traverse the singular type domains within one scope:

\begin{verbatim}
value = 42        # int
value = 3.14      # float
value = True      # bool
value = None      # NoneType
\end{verbatim}

Each assignment replaces the namespace entry. Python does not require separate names for separate roles, though separate names improve human readability when each binding represents a distinct pipeline stage (\texttt{raw\_score} -> \texttt{score} -> \texttt{passed}).

\subsubsection{Runtime Type Inspection with \texttt{type()} and \texttt{isinstance()}}

\texttt{type(x)} returns the exact type object of \texttt{x}. \texttt{isinstance(x, T)} asks whether \texttt{x} belongs to type \texttt{T} or a subclass.

The \texttt{bool}/\texttt{int} inheritance makes the distinction architecturally important:

\begin{verbatim}
flag = True
type(flag) is bool           # True
type(flag) is int            # False
isinstance(flag, bool)       # True
isinstance(flag, int)        # True  (subclass membership)
\end{verbatim}

Use \texttt{type()} for exact runtime type. Use \texttt{isinstance()} for category tests, but exclude \texttt{bool} explicitly when ``integer and not boolean'' is intended:

\begin{verbatim}
isinstance(x, (int, float)) and not isinstance(x, bool)
\end{verbatim}

The structural summary:

\begin{verbatim}
name = 42       -> references int object
name = "text"   -> same name, str object
type(name)      -> exact current type
isinstance(...) -> category including subclasses
\end{verbatim}

Python assignment is dynamic because the name can move. Python objects remain typed because each object carries its runtime type and behavior.
